\documentclass[runningheads]{llncs}
%
\usepackage{algpseudocode}
\usepackage{amsmath,algorithm,algorithmicx}
\usepackage{amsfonts,array}
\usepackage{textcomp}
\usepackage{caption,subcaption}
\usepackage{booktabs,multicol}
\usepackage{wrapfig,graphicx,pgfplots}
\usepackage{stfloats,minted}
\usepackage{cite,hyperref}
\usepackage{tikz,verbatim}
\usetikzlibrary{er,positioning,bayesnet}
\usepackage[inline]{enumitem} % Horizontal lists
\usepackage[all]{nowidow}
\usepackage[utf8]{vietnam}
\usepackage{lipsum} %just for generating dummy texts
%
\newcommand{\card}[1]{\left\vert{#1}\right\vert}
\newcommand*\Let[2]{\State #1 $\gets$ #2}
\definecolor{blue}{HTML}{1F77B4}
\definecolor{orange}{HTML}{FF7F0E}
\definecolor{green}{HTML}{2CA02C}
%
\pgfplotsset{compat=1.14}
%
\renewcommand{\topfraction}{0.85}
\renewcommand{\bottomfraction}{0.85}
\renewcommand{\textfraction}{0.15}
\renewcommand{\floatpagefraction}{0.8}
\renewcommand{\textfraction}{0.1}
\setlength{\floatsep}{3pt plus 1pt minus 1pt}
\setlength{\textfloatsep}{3pt plus 1pt minus 1pt}
\setlength{\intextsep}{3pt plus 1pt minus 1pt}
\setlength{\abovecaptionskip}{2pt plus 1pt minus 1pt}

%%%%%%%%%%%%%%%%%%%%%%%%%%%%%%%%%%%%%%%%%%%
\begin{document}
\title{Mẫu LaTeX Tuyển Tập KH\&CN Nhật Bản Vol.2}
%
\author{Nguyễn Văn A\inst{1} \and
Nguyễn Văn B\inst{2} \and
Nguyễn Văn C\inst{3}}
%
\institute{Đại Học Tokyo, Tokyo, Nhật Bản (\email{nva@utokyo.edu.jp})\\\and
Đại Học Osaka, Osaka, Nhật Bản (\email{nvb@uosaka.edu.jp})\\\and
Đại Học Điện Tử - Truyền Thông, Tokyo, Nhật Bản (\email{nvc@uec.ac.jp})}
%
\maketitle

%%%%%%%%%%%%%%%%%%%%%%%%%%%%%%%%%%%%%%%%%%%
\begin{abstract}
Phần abstract tóm tắt nội dung của bài viết trong khoảng 200-300 từ. \textbf{Bài viết mẫu này chỉ có tính chất tham khảo, các tác giả có thể điều chỉnh format và các đề mục phù hợp với nội dung của bài viết.} Nếu trong bài viết có sử dụng các hình ảnh, biểu đồ hoặc các thông tin khác của các bài báo hoặc sách đã xuất bản, tác giả phải tự xin phép các nhà xuất bản hoặc tác giả gốc theo quy định. VANJ sẽ không chịu trách nhiệm về các vấn đề nảy sinh liên quan đến bản quyền và có thể gỡ bài của tác giả nếu các quy định về bản quyền không được đảm bảo.

\keywords{\textbf{(3-6 từ)} vi mạch \and điện tử \and MEMS \and IoT \and quan trắc \and chăm sóc sức khỏe.}
\end{abstract}

%%%%%%%%%%%%%%%%%%%%%%%%%%%%%%%%%%%%%%%%%%%
\section{Giới Thiệu}
Một số lưu ý khi viết bài:
\begin{itemize}
    \item Các bài viết bằng tiếng việt cần đính kèm abstract bằng tiếng Anh (và cả tiếng Nhật nếu được).
    \item Các bài viết bằng tiếng Anh được khuyến khích đính kèm abstract bằng tiếng Việt, nếu không có sẽ được dịch thuật bởi HĐBT.
    \item Toàn văn bài viết vào khoảng 3000 đến 8000 từ.
    \item Chú ý: tối đa 2 hình vẽ hoặc bảng biểu trong 1 trang.
\end{itemize}

Ví dụ về buttlet point với số:
\begin{enumerate}
    \item The construction time.
    \item The accuracy of the cardinality estimates.
    \item The time needed to make a cardinality estimate.
    \item The number of values needed to store the model.
\end{enumerate}

Một ví dụ về bảng \ref{tab:1}:
\begin{table}[tbh]
\centering
\begin{tabular}{c|ccc}
& \textbf{Blond} & \textbf{Brown} & \textbf{Dark}\\\hline
\textbf{American} & 0.2 & 0.6 & 0.2\\\hline
\textbf{Swedish} & 0.8 & 0.2 & 0\\
\end{tabular}
\caption{$P(hair | nationality)$}
\label{tab:1}
\end{table}

Nếu như có nhiều hình. Các hình này cần được liệt kê và mô tả dưới dạng a), b) và c). Ví dự như trong hình \ref{fig:1}.

\begin{figure}[tbh]
    \centering
    \includegraphics[width=\textwidth]{figures/fig1.png}
    \caption{Hình vẽ trong bài viết  nên được đặt gần nơi được trích dẫn.}
    \label{fig:1}
\end{figure}

Công thức toán học được liệt kê và trích dẫn theo thứ tự trong bài viết như sau công thức \ref{eq:1}, công thức \ref{eq:2}.

\begin{equation}
\label{eq:1}
    Y = a+b+c
\end{equation}

\begin{equation}
\label{eq:2}
    X = a+b
\end{equation}

Vai trò, tầm quan trọng và triển vọng của công nghệ cảm biến \cite{ref1,ref2}.

Vị trí của Nhật Bản trong nghiên cứu về công nghệ cảm biến so với các nước trên thế giới.

Các thế mạnh của Nhật Bản trong công nghệ này

Tiềm năng áp dụng công nghệ về công nghệ cảm biến ở Việt Nam.

%%%%%%%%%%%%%%%%%%%%%%%%%%%%%%%%%%%%%%%%%%%
\section{Một Số Nghiên Cứu Tiêu Biểu của Nhật về Cảm Biến}

Nghiên cứu về nguyên lý của cảm biến \cite{ref3}.

Nguyên cứu về cách chế tạo cảm biến.

Nghiên cứu về ứng dụng của cảm biến.

%%%%%%%%%%%%%%%%%%%%%%%%%%%%%%%%%%%%%%%%%%%
\section{Một Số Lab Nghiên Cứu Tiêu Biểu của Nhật Bản Trong Lĩnh Vực Công Nghệ Cảm Biến}
\subsection{Lab Nghiên Cứu về Chế Tạo Cảm Biến tại Đại Học Tokyo}

Website: 

Giáo sư chủ nhiệm: 

Đề tài nghiên cứu chính: ...

%%%%%%%%%%%%%%%%%%%%%%%%%%%%%%%%%%%%%%%%%%%
\subsection{Lab Nghiên Cứu về Ứng Dụng Cảm Biến Trong Nông Nghiệp tại Đại Học Tokyo}

Website: 

Giáo sư chủ nhiệm: 

Đề tài nghiên cứu chính: ...

%%%%%%%%%%%%%%%%%%%%%%%%%%%%%%%%%%%%%%%%%%%
\section{Hướng Nghiên Cứu Tiềm Năng Ứng Dụng Công Nghệ Cảm Biến tại Việt Nam}
\subsection{Thực Trạng Nghiên Cứu Ứng Dụng Công Nghệ Cảm Biến tại Việt Nam Hiện Nay}
\lipsum[1-2]

%%%%%%%%%%%%%%%%%%%%%%%%%%%%%%%%%%%%%%%%%%%
\subsection{Đề Xuất các Phương Thức Áp Dụng Công Nghệ Cảm Biến tại Việt Nam}
\lipsum[3-4]

%%%%%%%%%%%%%%%%%%%%%%%%%%%%%%%%%%%%%%%%%%%
\section{Kết Luận}
\lipsum[5]

%%%%%%%%%%%%%%%%%%%%%%%%%%%%%%%%%%%%%%%%%%%
\section{Lời Cảm Ơn}
\lipsum[6]

%%%%%%%%%%%%%%%%%%%%%%%%%%%%%%%%%%%%%%%%%%%
\bibliography{myBib}{}
\bibliographystyle{IEEEtran}

\begin{wrapfigure}{l}{25mm} 
\includegraphics[width=1in,height=1.25in,clip,keepaspectratio]{figures/author.png}
\end{wrapfigure}
\par
\textbf{Nguyễn Văn A} (Nghiên cứu sinh) Mô tả ngắn thông tin của tác giả (Bio) trong khoảng 5-10 dòng. Link CV của tác giả hoặc các cách liên lạc khác như Google scholar, LinkedIn, Research gate, ORCID (sẽ hiện trên website Dự án và có thể được dùng trong các hoạt động truyền thông của Tuyển tập). Website của lab (nếu có)
\par

\end{document}
